\sectionright{\sffamily  Structural Elements}

%%%%%%%%%%%%%%%%%%%%%%%%%%%%%%%%%%%%%%%%%%%%%%%%%%%%%%%%%%%%%%%%%%%%%%%%%%%%%%%%%%%%%%%%%%%%%%%%%%%%%%%%%%%%%%%%%%%%%%%%%%%%%%%%%%%%%%%%%%%%%%%%%%%%%%%%%%%%%%%%%%%%%%%%%
\subsection{\sffamily Custom Module (Black Box)}

    \subsubsection{\sffamily Description}
        The \texttt{rtl\_module} element is designed to represent higher-level functional units or "black boxes" within a digital system. It is based on the \texttt{rtl\_base\_box} but features a distinct green tint and a dedicated label at the top. 
        This element is ideal for hierarchical diagrams where the internal logic of a component is abstracted away. The label is automatically positioned at the \texttt{north} edge with a slight offset to prevent overlap with incoming signals.

    \subsubsection{\sffamily Attributes}
        The module's name is passed as a mandatory argument to the style. Dimensioning follows the standard package scheme:

    \begin{center}
        \begin{tabular}{lp{8cm}}
        \toprule
            \textbf{Attribute} & \textbf{Effect} \\
            \midrule
            \texttt{rtl\_module=\{Name\}} & Sets the identifying label at the top of the module. \\
            \texttt{rtl\_width} & Sets the minimum width (Default: 1.5cm). \\
            \texttt{rtl\_height} & Sets the minimum height (Default: 1.8cm). \\
            \bottomrule
        \end{tabular}
    \end{center}

    \subsubsection{\sffamily Anchors}
        \begin{itemize}
            \item \textbf{Standard Anchors}: \texttt{north}, \texttt{south}, \texttt{east}, \texttt{west}, and all corner anchors.
            \item \textbf{Note}: Use the \texttt{north} anchor with caution, as it shares space with the module label.
        \end{itemize}

    \subsubsection{\sffamily Listing}
        The following example demonstrates a scaled module with a custom name.

\begin{RTLexample}[Custom Module]
\begin{tikzpicture}[>=Stealth, thick, scale=0.8]
    % Module with specific name and dimensions
    \node[rtl_module={Control\_Unit}, rtl_width=3cm, rtl_height=2.5cm] (Mod1) at (0,0) {};
    
    % Connections
    \draw[<-] (Mod1.west) ++(0,0.5) -- ++(-1,0) node[left] {instr\_in};
    \draw[->] (Mod1.east) -- ++(1,0) node[right] {ctrl\_out};
    \draw[<-] (Mod1.south) -- ++(0,-0.7) node[below] {reset};
\end{tikzpicture}
\end{RTLexample}

%%%%%%%%%%%%%%%%%%%%%%%%%%%%%%%%%%%%%%%%%%%%%%%%%%%%%%%%%%%%%%%%%%%%%%%%%%%%%%%%%%%%%%%%%%%%%%%%%%%%%%%%%%%%%%%%%%%%%%%%%%%%%%%%%%%%%%%%%%%%%%%%%%%%%%%%%%%%%%%%%%%%%%%%%
\subsection{\sffamily Constant Value}

    \subsubsection{\sffamily Description}
        The \texttt{rtl\_const} element is used to represent static signal values, parameters, or hardcoded addresses. Its distinct purple color serves as a visual cue that the signal originating from this block is non-volatile and fixed during the design's operation. 
        Internally, it behaves like an asynchronous register without any clock or control logic.

    \subsubsection{\sffamily Attributes}
    \begin{center}
        \begin{tabular}{lp{8cm}}
        \toprule
            \textbf{Attribute} & \textbf{Effect} \\
            \midrule
            \texttt{rtl\_width} & Sets the minimum width (Default: 1.5cm). \\
            \texttt{rtl\_height} & Sets the minimum height (Default: 1.8cm). \\
            \bottomrule
        \end{tabular}
    \end{center}

    \subsubsection{\sffamily Anchors}
        \begin{itemize}
            \item \textbf{Standard Anchors}: Full set of rectangular anchors (\texttt{north}, \texttt{south}, etc.).
        \end{itemize}

    \subsubsection{\sffamily Listing}
        Constants are typically used to feed fixed values into operators or modules.

\begin{RTLexample}[Constant Value]
\begin{tikzpicture}[>=Stealth, thick, scale=0.8]
    % Constant value block
    \node[rtl_const] (C1) at (0,0) {8'hFF};
    
    % Example of signal origin
    \draw[->, rtl_bus={8}] (C1.east) -- ++(1.5,0) node[right] {};
    
\end{tikzpicture}
\end{RTLexample}

%%%%%%%%%%%%%%%%%%%%%%%%%%%%%%%%%%%%%%%%%%%%%%%%%%%%%%%%%%%%%%%%%%%%%%%%%%%%%%%%%%%%%%%%%%%%%%%%%%%%%%%%%%%%%%%%%%%%%%%%%%%%%%%%%%%%%%%%%%%%%%%%%%%%%%%%%%%%%%%%%%%%%%%%%
\subsection{\sffamily Cutset Decoration}

    \subsubsection{\sffamily Description}
        The \texttt{is\_cutset} style is a specialized visual decorator used to mark timing boundaries, pipelining stages, or logic cuts within a diagram. Instead of being a standalone node, it is applied as an attribute to any existing rectangular element (like registers or modules).
        
        The style automatically draws a dashed vertical line through the center of the component, extending slightly (5\%) beyond the \texttt{north} and \texttt{south} boundaries. By utilizing the \texttt{foreground} PGF layer, the cutset remains clearly visible even when applied to filled or shaded elements.

    \subsubsection{\sffamily Attributes}
        The style accepts an optional argument to customize the appearance of the dashed line, such as changing its color or thickness.

    \begin{center}
        \begin{tabular}{lp{8cm}}
        \toprule
            \textbf{Attribute} & \textbf{Effect} \\
            \midrule
            \texttt{is\_cutset=\{draw\_options\}} & Activates the cutset line. The argument allows standard TikZ draw options (Default: \texttt{dashed, thick}). \\
            \bottomrule
        \end{tabular}
    \end{center}

    \subsubsection{\sffamily Anchors}
            \begin{itemize}
                \item \textbf{Standard Anchors}: Full set of rectangular anchors (\texttt{north}, \texttt{south}, etc.).
            \end{itemize}

    \subsubsection{\sffamily Listing}
        The following example shows how to mark a timing boundary across a series of registers.

\begin{RTLexample}[Logic Cutset]
\begin{tikzpicture}[>=Stealth, thick, scale=0.8]
    % Sequence of registers
    \node[rtl_module={Mod A}] (M1) at (0,0) {};
    
    % Register with a red cutset line
    \node[rtl_reg, is_cutset={red}, right=2cm of M1] (R1) {Reg};
    
    \node[rtl_module={Mod B}, right=2cm of R1] (M2) {};

    % Signals
    \draw[->] (M1) -- (R1);
    \draw[->] (R1) -- (M2);

    % Label for the cutset
    \node[red, font=\sffamily\small] at (R1.north) [yshift=0.5cm] {Timing Cut};
\end{tikzpicture}
\end{RTLexample}
%%%%%%%%%%%%%%%%%%%%%%%%%%%%%%%%%%%%%%%%%%%%%%%%%%%%%%%%%%%%%%%%%%%%%%%%%%%%%%%%%%%%%%%%%%%%%%%%%%%%%%%%%%%%%%%%%%%%%%%%%%%%%%%%%%%%%%%%%%%%%%%%%%%%%%%%%%%%%%%%%%%%%%%%%

%%%%%%%%%%%%%%%%%%%%%%%%%%%%%%%%%%%%%%%%%%%%%%%%%%%%%%%%%%%%%%%%%%%%%%%%%%%%%%%%%%%%%%%%%%%%%%%%%%%%%%%%%%%%%%%%%%%%%%%%%%%%%%%%%%%%%%%%%%%%%%%%%%%%%%%%%%%%%%%%%%%%%%%%%
\subsection{\sffamily Hierarchical Levels}

    \subsubsection{\sffamily Description}
        Complex digital systems, such as SOC architectures or multi-stage DSP pipelines, etc., often require a hierarchical representation to remain readable. The \texttt{\textbackslash createHierLevel} command is designed to automate the "zoom-in" effect from a high-level block to a detailed schematic. 
        
        It automatically draws a focal marker (dashed circle) around the source component, creates a large container at a target position, and pulls projection lines on the \texttt{background} layer to visualize the transition. This command is specifically optimized to include external PDF schematics, allowing for a modular documentation workflow.

    \subsubsection{\sffamily Attributes}
        The command utilizes a mix of positional and mandatory arguments to define the geometric relationship between the hierarchy levels:

    \begin{center}
        \begin{tabular}{lp{8cm}}
        \toprule
            \textbf{Argument} & \textbf{Effect} \\
            \midrule
            \texttt{[\#1]} & Optional TikZ draw options for the projection lines (e.g., color, dash pattern). \\
            \texttt{\{\#2\}} & Name of the source node in the overview (e.g., \texttt{DSP\_Core}). \\
            \texttt{\{\#3\}} & Target coordinate for the center of the detail view (e.g., \texttt{(10,-5)}). \\
            \texttt{\{\#4\}} & Reference name for the new detail container. \\
            \texttt{\{\#5\}} & Width of the detail container. \\
            \texttt{\{\#6\}} & Height of the detail container. \\
            \texttt{\{\#7\}} & File path to the external PDF or image. \\
            \bottomrule
        \end{tabular}
    \end{center}

    \subsubsection{\sffamily Anchors}
        The created detail container behaves like a standard \texttt{rtl\_reg} node:
        \begin{itemize}
            \item \textbf{Standard Anchors}: \texttt{north}, \texttt{south}, \texttt{east}, \texttt{west}, etc. These can be used to connect global signals to the zoomed-in schematic.
        \end{itemize}

    \subsubsection{\sffamily Listing}
        The following example demonstrates a zoom-in from a generic DSP block to a specific 8-to-1 Multiplexer implementation.

\begin{RTLexample}[Hierarchical Level]
\begin{tikzpicture}[>=Stealth, thick, scale=0.8]
    \node[rtl_module] (Core) at (0,0) {Module};

    \createHierLevel{Core}{(12,-4)}{DetailView}{8cm}{7cm}{../examples/8_1_MUX/8_1_MUX.pdf};

\end{tikzpicture}
\end{RTLexample}